\documentclass[aspectratio=169, 11pt]{beamer}
% ============================================================================
% Preambulo compartido para presentaciones Beamer
% Estadistica 2026-II - Pontificia Universidad Javeriana
% Incluir con: % ============================================================================
% Preambulo compartido para presentaciones Beamer
% Estadistica 2026-II - Pontificia Universidad Javeriana
% Incluir con: % ============================================================================
% Preambulo compartido para presentaciones Beamer
% Estadistica 2026-II - Pontificia Universidad Javeriana
% Incluir con: \input{../preambulo_beamer.tex} despues de \documentclass
% ============================================================================

\usepackage[T1]{fontenc}
\usepackage[utf8]{inputenc}
\usepackage[spanish]{babel}
\usepackage{amsmath, amssymb, amsfonts}
\usepackage{booktabs}
\usepackage{graphicx}
\usepackage{tikz}
\usetikzlibrary{shapes.geometric, positioning, arrows.meta, calc}
\usepackage{pgfplots}
\pgfplotsset{compat=1.18}
\usepackage{listings}
\usepackage{xcolor}
\usepackage{hyperref}
\usepackage{multicol}

% --- Tema Beamer (minimalista) ---
\usetheme{default}
\usecolortheme{default}
\usefonttheme{default}
\setbeamertemplate{navigation symbols}{}
\setbeamertemplate{caption}[numbered]
\setbeamertemplate{footline}{%
  \hspace{1em}{\tiny\url{https://github.com/polanco-jaime/Curso-Estadistica-2026-3}}\hfill\usebeamerfont{footline}\insertframenumber/\inserttotalframenumber\hspace{1em}\vspace{0.5em}%
}
\setbeamertemplate{itemize items}[circle]
\setbeamertemplate{enumerate items}[default]

% --- Colores institucionales (sutiles) ---
\definecolor{javeriana}{RGB}{0, 62, 105}
\definecolor{javerianalight}{RGB}{0, 105, 170}
\definecolor{codigobg}{RGB}{245, 245, 245}
\definecolor{codigofg}{RGB}{40, 40, 40}
\definecolor{comentario}{RGB}{100, 100, 100}
\definecolor{keyword}{RGB}{0, 100, 150}
\definecolor{string}{RGB}{180, 60, 30}

\setbeamercolor{title}{fg=javeriana}
\setbeamercolor{frametitle}{fg=javeriana}
\setbeamercolor{structure}{fg=javeriana}
\setbeamercolor{block title}{fg=white, bg=javeriana!75}
\setbeamercolor{block body}{bg=javeriana!5}
\setbeamercolor{block title alerted}{fg=white, bg=red!70!black}
\setbeamercolor{block body alerted}{bg=red!5}
\setbeamercolor{block title example}{fg=white, bg=green!40!black}
\setbeamercolor{block body example}{bg=green!5}
\setbeamercolor{item}{fg=javeriana}

\setbeamerfont{title}{series=\bfseries, size=\Large}
\setbeamerfont{frametitle}{series=\bfseries}

% --- Configuracion de listings para R ---
\lstset{
  language=R,
  basicstyle=\ttfamily\footnotesize,
  backgroundcolor=\color{codigobg},
  commentstyle=\color{comentario}\itshape,
  keywordstyle=\color{keyword}\bfseries,
  stringstyle=\color{string},
  numbers=none,
  frame=single,
  rulecolor=\color{javeriana!30},
  breaklines=true,
  showstringspaces=false,
  columns=flexible,
  morekeywords={library, ggplot, aes, geom_histogram, geom_boxplot,
    geom_point, geom_smooth, geom_density, geom_bar, geom_line,
    geom_vline, geom_hline, geom_errorbar, labs, theme_minimal,
    facet_wrap, scale_fill_manual, scale_color_manual,
    summarise, group_by, filter, mutate, select, arrange,
    read.csv, head, str, summary, mean, median, sd, var, cor,
    lm, t.test, chisq.test, wilcox.test, kruskal.test, aov,
    pnorm, qnorm, dnorm, rnorm, qt, pt, qchisq, pchisq, qf, pf,
    confint, coeftest, vcovHC, bptest, vif, cooks.distance,
    TukeyHSD, table, prop.table}
}

% --- Comandos personalizados ---
\newcommand{\R}{\texttt{R}}
\newcommand{\code}[1]{\texttt{\footnotesize #1}}
\newcommand{\variable}[1]{\texttt{#1}}
\newcommand{\paquete}[1]{\texttt{#1}}

% --- Informacion del curso ---
\institute{Pontificia Universidad Javeriana\\
  Facultad de Ciencias Econ\'omicas y Administrativas}
\date{2026-II}
 despues de \documentclass
% ============================================================================

\usepackage[T1]{fontenc}
\usepackage[utf8]{inputenc}
\usepackage[spanish]{babel}
\usepackage{amsmath, amssymb, amsfonts}
\usepackage{booktabs}
\usepackage{graphicx}
\usepackage{tikz}
\usetikzlibrary{shapes.geometric, positioning, arrows.meta, calc}
\usepackage{pgfplots}
\pgfplotsset{compat=1.18}
\usepackage{listings}
\usepackage{xcolor}
\usepackage{hyperref}
\usepackage{multicol}

% --- Tema Beamer (minimalista) ---
\usetheme{default}
\usecolortheme{default}
\usefonttheme{default}
\setbeamertemplate{navigation symbols}{}
\setbeamertemplate{caption}[numbered]
\setbeamertemplate{footline}{%
  \hspace{1em}{\tiny\url{https://github.com/polanco-jaime/Curso-Estadistica-2026-3}}\hfill\usebeamerfont{footline}\insertframenumber/\inserttotalframenumber\hspace{1em}\vspace{0.5em}%
}
\setbeamertemplate{itemize items}[circle]
\setbeamertemplate{enumerate items}[default]

% --- Colores institucionales (sutiles) ---
\definecolor{javeriana}{RGB}{0, 62, 105}
\definecolor{javerianalight}{RGB}{0, 105, 170}
\definecolor{codigobg}{RGB}{245, 245, 245}
\definecolor{codigofg}{RGB}{40, 40, 40}
\definecolor{comentario}{RGB}{100, 100, 100}
\definecolor{keyword}{RGB}{0, 100, 150}
\definecolor{string}{RGB}{180, 60, 30}

\setbeamercolor{title}{fg=javeriana}
\setbeamercolor{frametitle}{fg=javeriana}
\setbeamercolor{structure}{fg=javeriana}
\setbeamercolor{block title}{fg=white, bg=javeriana!75}
\setbeamercolor{block body}{bg=javeriana!5}
\setbeamercolor{block title alerted}{fg=white, bg=red!70!black}
\setbeamercolor{block body alerted}{bg=red!5}
\setbeamercolor{block title example}{fg=white, bg=green!40!black}
\setbeamercolor{block body example}{bg=green!5}
\setbeamercolor{item}{fg=javeriana}

\setbeamerfont{title}{series=\bfseries, size=\Large}
\setbeamerfont{frametitle}{series=\bfseries}

% --- Configuracion de listings para R ---
\lstset{
  language=R,
  basicstyle=\ttfamily\footnotesize,
  backgroundcolor=\color{codigobg},
  commentstyle=\color{comentario}\itshape,
  keywordstyle=\color{keyword}\bfseries,
  stringstyle=\color{string},
  numbers=none,
  frame=single,
  rulecolor=\color{javeriana!30},
  breaklines=true,
  showstringspaces=false,
  columns=flexible,
  morekeywords={library, ggplot, aes, geom_histogram, geom_boxplot,
    geom_point, geom_smooth, geom_density, geom_bar, geom_line,
    geom_vline, geom_hline, geom_errorbar, labs, theme_minimal,
    facet_wrap, scale_fill_manual, scale_color_manual,
    summarise, group_by, filter, mutate, select, arrange,
    read.csv, head, str, summary, mean, median, sd, var, cor,
    lm, t.test, chisq.test, wilcox.test, kruskal.test, aov,
    pnorm, qnorm, dnorm, rnorm, qt, pt, qchisq, pchisq, qf, pf,
    confint, coeftest, vcovHC, bptest, vif, cooks.distance,
    TukeyHSD, table, prop.table}
}

% --- Comandos personalizados ---
\newcommand{\R}{\texttt{R}}
\newcommand{\code}[1]{\texttt{\footnotesize #1}}
\newcommand{\variable}[1]{\texttt{#1}}
\newcommand{\paquete}[1]{\texttt{#1}}

% --- Informacion del curso ---
\institute{Pontificia Universidad Javeriana\\
  Facultad de Ciencias Econ\'omicas y Administrativas}
\date{2026-II}
 despues de \documentclass
% ============================================================================

\usepackage[T1]{fontenc}
\usepackage[utf8]{inputenc}
\usepackage[spanish]{babel}
\usepackage{amsmath, amssymb, amsfonts}
\usepackage{booktabs}
\usepackage{graphicx}
\usepackage{tikz}
\usetikzlibrary{shapes.geometric, positioning, arrows.meta, calc}
\usepackage{pgfplots}
\pgfplotsset{compat=1.18}
\usepackage{listings}
\usepackage{xcolor}
\usepackage{hyperref}
\usepackage{multicol}

% --- Tema Beamer (minimalista) ---
\usetheme{default}
\usecolortheme{default}
\usefonttheme{default}
\setbeamertemplate{navigation symbols}{}
\setbeamertemplate{caption}[numbered]
\setbeamertemplate{footline}{%
  \hspace{1em}{\tiny\url{https://github.com/polanco-jaime/Curso-Estadistica-2026-3}}\hfill\usebeamerfont{footline}\insertframenumber/\inserttotalframenumber\hspace{1em}\vspace{0.5em}%
}
\setbeamertemplate{itemize items}[circle]
\setbeamertemplate{enumerate items}[default]

% --- Colores institucionales (sutiles) ---
\definecolor{javeriana}{RGB}{0, 62, 105}
\definecolor{javerianalight}{RGB}{0, 105, 170}
\definecolor{codigobg}{RGB}{245, 245, 245}
\definecolor{codigofg}{RGB}{40, 40, 40}
\definecolor{comentario}{RGB}{100, 100, 100}
\definecolor{keyword}{RGB}{0, 100, 150}
\definecolor{string}{RGB}{180, 60, 30}

\setbeamercolor{title}{fg=javeriana}
\setbeamercolor{frametitle}{fg=javeriana}
\setbeamercolor{structure}{fg=javeriana}
\setbeamercolor{block title}{fg=white, bg=javeriana!75}
\setbeamercolor{block body}{bg=javeriana!5}
\setbeamercolor{block title alerted}{fg=white, bg=red!70!black}
\setbeamercolor{block body alerted}{bg=red!5}
\setbeamercolor{block title example}{fg=white, bg=green!40!black}
\setbeamercolor{block body example}{bg=green!5}
\setbeamercolor{item}{fg=javeriana}

\setbeamerfont{title}{series=\bfseries, size=\Large}
\setbeamerfont{frametitle}{series=\bfseries}

% --- Configuracion de listings para R ---
\lstset{
  language=R,
  basicstyle=\ttfamily\footnotesize,
  backgroundcolor=\color{codigobg},
  commentstyle=\color{comentario}\itshape,
  keywordstyle=\color{keyword}\bfseries,
  stringstyle=\color{string},
  numbers=none,
  frame=single,
  rulecolor=\color{javeriana!30},
  breaklines=true,
  showstringspaces=false,
  columns=flexible,
  morekeywords={library, ggplot, aes, geom_histogram, geom_boxplot,
    geom_point, geom_smooth, geom_density, geom_bar, geom_line,
    geom_vline, geom_hline, geom_errorbar, labs, theme_minimal,
    facet_wrap, scale_fill_manual, scale_color_manual,
    summarise, group_by, filter, mutate, select, arrange,
    read.csv, head, str, summary, mean, median, sd, var, cor,
    lm, t.test, chisq.test, wilcox.test, kruskal.test, aov,
    pnorm, qnorm, dnorm, rnorm, qt, pt, qchisq, pchisq, qf, pf,
    confint, coeftest, vcovHC, bptest, vif, cooks.distance,
    TukeyHSD, table, prop.table}
}

% --- Comandos personalizados ---
\newcommand{\R}{\texttt{R}}
\newcommand{\code}[1]{\texttt{\footnotesize #1}}
\newcommand{\variable}[1]{\texttt{#1}}
\newcommand{\paquete}[1]{\texttt{#1}}

% --- Informacion del curso ---
\institute{Pontificia Universidad Javeriana\\
  Facultad de Ciencias Econ\'omicas y Administrativas}
\date{2026-II}


\title{Sesi\'on 0: Conceptos Fundamentales}
\subtitle{Lenguaje b\'asico de probabilidad y variables aleatorias}
\author{Prof.\ Jaime Polanco Jim\'enez}
\date{2026-II}

\begin{document}

\begin{frame}
\titlepage
\end{frame}

\begin{frame}{Agenda}
\tableofcontents
\end{frame}

% ============================================================================
\section{Experimento aleatorio}
% ============================================================================

\begin{frame}{Experimento aleatorio}
\begin{columns}[T]
\column{0.48\textwidth}
\begin{block}{Definici\'on formal}
Un \alert{experimento aleatorio} es un proceso cuyo resultado no puede predecirse con certeza antes de realizarlo.

\vspace{0.5em}
Propiedades:
\begin{enumerate}
  \item Se conocen todos los resultados posibles.
  \item No se sabe cu\'al ocurrir\'a.
  \item Puede repetirse bajo las mismas condiciones.
\end{enumerate}
\end{block}

\column{0.48\textwidth}
\begin{exampleblock}{Ejemplo: subir un reel a Instagram}
Subes un reel a Instagram.

\vspace{0.5em}
\begin{itemize}
  \item Sabes que puede tener \textbf{muchas} o \textbf{pocas} views.
  \item Pero al publicarlo, \textbf{no sabes} cu\'antas tendr\'a.
  \item Si subes otro reel ma\~nana, el resultado puede ser distinto.
\end{itemize}

\vspace{0.5em}
\textit{Antes de presionar ``compartir'', el resultado es incierto.}
\end{exampleblock}
\end{columns}
\end{frame}

% ============================================================================
\section{Espacio muestral}
% ============================================================================

\begin{frame}{Espacio muestral $\Omega$}
\begin{columns}[T]
\column{0.48\textwidth}
\begin{block}{Definici\'on formal}
El \alert{espacio muestral} $\Omega$ es el conjunto de \textbf{todos} los resultados posibles de un experimento aleatorio.

\vspace{0.5em}
Puede ser:
\begin{itemize}
  \item \textbf{Finito:} $\Omega = \{\omega_1, \omega_2, \dots, \omega_n\}$
  \item \textbf{Numerable:} $\Omega = \{\omega_1, \omega_2, \dots\}$
  \item \textbf{No numerable:} $\Omega \subseteq \mathbb{R}$
\end{itemize}
\end{block}

\column{0.48\textwidth}
\begin{exampleblock}{Ejemplo: views de un reel}
Si subimos un reel, las views posibles son:
\[
  \Omega = \{0, 1, 2, 3, \dots\}
\]

\vspace{0.3em}
Si solo nos interesa si \textbf{se vuelve viral}:
\[
  \Omega = \{\text{viral},\; \text{no viral}\}
\]

\vspace{0.3em}
\textit{$\Omega$ delimita el ``universo'' del experimento.}
\end{exampleblock}
\end{columns}
\end{frame}

% ============================================================================
\section{Evento}
% ============================================================================

\begin{frame}{Evento}
\begin{columns}[T]
\column{0.48\textwidth}
\begin{block}{Definici\'on formal}
Un \alert{evento} $A$ es un subconjunto del espacio muestral:
\[
  A \subseteq \Omega
\]

\vspace{0.3em}
La colecci\'on de eventos forma una $\sigma$-\'algebra $\mathcal{A}$ sobre $\Omega$:
\begin{enumerate}
  \item $\Omega \in \mathcal{A}$
  \item Si $A \in \mathcal{A}$, entonces $A^c \in \mathcal{A}$
  \item Si $A_1, A_2, \dots \in \mathcal{A}$, entonces $\bigcup_{i=1}^{\infty} A_i \in \mathcal{A}$
\end{enumerate}
\end{block}

\column{0.48\textwidth}
\begin{exampleblock}{Ejemplo: views de un reel}
Evento $A$: ``que el reel supere 10{,}000 views''.
\[
  A = \{\omega \in \Omega : \omega > 10{,}000\}
\]

\vspace{0.3em}
Evento $A^c$: ``que se quede en 10{,}000 views o menos''.

\vspace{0.3em}
\textit{Un evento es una pregunta de s\'i/no sobre el resultado.}
\end{exampleblock}
\end{columns}
\end{frame}

% ============================================================================
\section{Probabilidad}
% ============================================================================

\begin{frame}{Probabilidad $P$}
\begin{columns}[T]
\column{0.48\textwidth}
\begin{block}{Axiomas de Kolmogorov}
Una \alert{medida de probabilidad} es una funci\'on $P: \mathcal{A} \to [0,1]$ tal que:

\vspace{0.3em}
\begin{enumerate}
  \item $P(A) \geq 0$ para todo $A \in \mathcal{A}$
  \item $P(\Omega) = 1$
  \item Si $A_1, A_2, \dots$ son disjuntos:
  \[
    P\!\left(\bigcup_{i=1}^{\infty} A_i\right) = \sum_{i=1}^{\infty} P(A_i)
  \]
\end{enumerate}

\vspace{0.3em}
La terna $(\Omega, \mathcal{A}, P)$ es un \alert{espacio de probabilidad}.
\end{block}

\column{0.48\textwidth}
\begin{exampleblock}{Ejemplo: probabilidad de hacerse viral}
?`Qu\'e tan probable es que tu reel se vuelva viral?

\vspace{0.3em}
Basado en tu historial de reels:
\begin{itemize}
  \item $P(\text{viral}) = 0.05$
  \item $P(\text{no viral}) = 0.95$
  \item Suman 1 \checkmark
\end{itemize}

\vspace{0.3em}
$P$ asigna un \textbf{n\'umero entre 0 y 1} a cada evento, reflejando qu\'e tan plausible es.

\vspace{0.3em}
\textit{0 = imposible, 1 = seguro.}
\end{exampleblock}
\end{columns}
\end{frame}

% ============================================================================
\section{Variable aleatoria}
% ============================================================================

\begin{frame}{Variable aleatoria}
\begin{columns}[T]
\column{0.48\textwidth}
\begin{block}{Definici\'on formal}
Una \alert{variable aleatoria} sobre $(\Omega, \mathcal{A}, P)$ es una funci\'on medible:
\[
  X: (\Omega, \mathcal{A}) \to (\mathbb{R}, \mathcal{B}(\mathbb{R}))
\]

tal que:
\[
  X^{-1}(B) \in \mathcal{A}, \quad \forall\, B \in \mathcal{B}(\mathbb{R})
\]

donde $\mathcal{B}(\mathbb{R})$ es la $\sigma$-\'algebra de Borel.

\vspace{0.3em}
\textit{$X$ traduce resultados del experimento a n\'umeros reales.}
\end{block}

\column{0.48\textwidth}
\begin{exampleblock}{Ejemplo: likes en Instagram}
Definimos $X$ = n\'umero de likes en una foto.
\[
  X: \Omega \to \{0, 1, 2, \dots\}
\]

\vspace{0.3em}
\begin{itemize}
  \item El resultado era \textit{cualitativo} (``la gente reaccion\'o'').
  \item $X$ lo convierte en un \textbf{n\'umero}: 47, 230, 1{,}500\dots
  \item Ahora podemos hacer \textbf{matem\'aticas} con \'el.
\end{itemize}

\vspace{0.3em}
\textit{La variable aleatoria es el puente entre el fen\'omeno y los n\'umeros.}
\end{exampleblock}
\end{columns}
\end{frame}

% ============================================================================
\section{Tipos de variable aleatoria}
% ============================================================================

\begin{frame}{Tipos de variable aleatoria}
\begin{columns}[T]
\column{0.48\textwidth}
\begin{block}{Discreta vs.\ Continua}
\textbf{VA Discreta:} toma valores en un conjunto finito o numerable.
\[
  X: \Omega \to \{x_1, x_2, \dots\}
\]

\vspace{0.3em}
\textbf{VA Continua:} toma valores en un intervalo de $\mathbb{R}$.
\[
  X: \Omega \to \mathbb{R} \quad (\text{rango no numerable})
\]

\vspace{0.3em}
La diferencia determina si usamos \textbf{sumas} o \textbf{integrales}.
\end{block}

\column{0.48\textwidth}
\begin{exampleblock}{Ejemplos cotidianos}
\textbf{Discreta:}
\begin{itemize}
  \item N\'umero de matches en una app de citas: 0, 1, 2, 3, \dots
  \item Cantidad de kills en una partida de videojuegos.
\end{itemize}

\vspace{0.5em}
\textbf{Continua:}
\begin{itemize}
  \item Horas que pasas hoy en TikTok: 0.5, 1.3, 2.7\dots
  \item Tiempo de entrega de un pedido por app (Rappi, etc.).
\end{itemize}

\vspace{0.3em}
\textit{?`Se puede contar? $\to$ discreta. ?`Se mide? $\to$ continua.}
\end{exampleblock}
\end{columns}
\end{frame}

% ============================================================================
\section{Distribuci\'on de probabilidad}
% ============================================================================

\begin{frame}{Distribuci\'on de probabilidad}
\begin{columns}[T]
\column{0.48\textwidth}
\begin{block}{Definici\'on formal}
\textbf{Discreta} --- funci\'on de masa (pmf):
\[
  p(x) = P(X = x), \quad \sum_{x} p(x) = 1
\]
\textbf{Continua} --- funci\'on de densidad (pdf):
\[
  f(x) \geq 0, \quad \int_{-\infty}^{\infty}\! f(x)\,dx = 1
\]
\[
  P(a \leq X \leq b) = \int_a^b f(x)\,dx
\]
\end{block}

\column{0.48\textwidth}
\begin{exampleblock}{Ejemplo: likes y screen time}
\textbf{Discreta (likes en una foto):}

\vspace{0.2em}
{\small\begin{tabular}{cc}
\toprule
$x$ (likes) & $p(x)$ \\
\midrule
150 & 0.40 \\
200 & 0.30 \\
300 & 0.20 \\
500+ & 0.10 \\
\bottomrule
\end{tabular}}

\vspace{0.3em}
\textbf{Continua:} densidad del screen time diario (horas en el celular).

\vspace{0.2em}
\textit{La distribuci\'on describe el ``patr\'on'' de los resultados.}
\end{exampleblock}
\end{columns}
\end{frame}

% ============================================================================
\section{Valor esperado y varianza}
% ============================================================================

\begin{frame}{Valor esperado $E[X]$}
\begin{columns}[T]
\column{0.48\textwidth}
\begin{block}{Definici\'on formal}
\textbf{Discreta:}
\[
  E[X] = \sum_{x} x \cdot p(x)
\]

\textbf{Continua:}
\[
  E[X] = \int_{-\infty}^{\infty} x \cdot f(x)\,dx
\]

\vspace{0.3em}
$E[X]$ es el \alert{promedio ponderado} de los valores que $X$ puede tomar, ponderados por su probabilidad.
\end{block}

\column{0.48\textwidth}
\begin{exampleblock}{Ejemplo: Spotify Wrapped}
Tu Spotify Wrapped dice que escuchas en promedio 85 min de m\'usica al d\'ia.

\vspace{0.3em}
Si $X$ = minutos diarios de m\'usica:
\[
  E[X] = 85 \text{ min}
\]

\vspace{0.3em}
Interpretaci\'on: si sumas todos los d\'ias del a\~no y divides entre 365, obtienes $\sim$85 min.

\vspace{0.5em}
\textit{$E[X]$ es el ``centro de gravedad'' de la distribuci\'on.}
\end{exampleblock}
\end{columns}
\end{frame}

\begin{frame}{Varianza $\text{Var}(X)$}
\begin{columns}[T]
\column{0.48\textwidth}
\begin{block}{Definici\'on formal}
\[
  \text{Var}(X) = E\!\left[(X - E[X])^2\right]
\]

Equivalentemente:
\[
  \text{Var}(X) = E[X^2] - (E[X])^2
\]

\vspace{0.3em}
La \alert{desviaci\'on est\'andar}:
\[
  \sigma = \sqrt{\text{Var}(X)}
\]

mide la dispersi\'on en las mismas unidades que $X$.
\end{block}

\column{0.48\textwidth}
\begin{exampleblock}{Ejemplo: Spotify Wrapped}
D\'ias estables: escuchas $\sim$60 min $\to$ varianza \textbf{baja}.

\vspace{0.3em}
D\'ias de \textit{binge listening}: 20 min un d\'ia, 200 min otro $\to$ varianza \textbf{alta}.

\vspace{0.3em}
Dos personas pueden tener $E[X] = 85$ min, pero una es \textbf{constante} y la otra \textbf{impredecible}.

\vspace{0.5em}
\textit{La varianza mide qu\'e tan \textbf{impredecible} es el resultado.}
\end{exampleblock}
\end{columns}
\end{frame}

% ============================================================================
\section{Resumen}
% ============================================================================

\begin{frame}{Resumen: del experimento al an\'alisis}
\begin{center}
\begin{tikzpicture}[
  box/.style={rectangle, draw=javeriana, fill=javeriana!8, rounded corners,
              minimum width=2.8cm, minimum height=0.9cm, align=center,
              font=\small},
  arr/.style={-{Stealth[length=3mm]}, thick, javeriana}
]
\node[box] (exp) at (0,0) {Experimento\\aleatorio};
\node[box] (omega) at (3.5,0) {Espacio muestral\\$\Omega$};
\node[box] (evento) at (7,0) {Evento\\$A \subseteq \Omega$};
\node[box] (prob) at (10.5,0) {Probabilidad\\$P(A)$};
\node[box] (va) at (3.5,-2) {Variable aleatoria\\$X: \Omega \to \mathbb{R}$};
\node[box] (dist) at (7,-2) {Distribuci\'on\\pmf / pdf};
\node[box] (mom) at (10.5,-2) {$E[X]$, $\text{Var}(X)$};

\draw[arr] (exp) -- (omega);
\draw[arr] (omega) -- (evento);
\draw[arr] (evento) -- (prob);
\draw[arr] (omega) -- (va);
\draw[arr] (va) -- (dist);
\draw[arr] (dist) -- (mom);
\end{tikzpicture}
\end{center}

\vspace{0.5em}
\begin{itemize}
  \item Estos conceptos son el \textbf{lenguaje} que usaremos todo el semestre.
  \item A partir de la \alert{Sesi\'on 1} aplicaremos estas ideas a \textbf{datos reales}.
\end{itemize}
\end{frame}

\begin{frame}{}
\vfill
\begin{center}
{\Large\color{javeriana} ?`Preguntas?}

\vspace{1em}
{\normalsize Siguiente: \textbf{Sesi\'on 1 --- Estad\'istica Descriptiva Num\'erica}}
\end{center}
\vfill
\end{frame}

\end{document}
